%%==================================%%
%%           第 7 章  结论           %%
%%==================================%%
\chapter{总结与展望}

\section{系统总结}

\subsection{主要功能完成情况}

经过较为系统的设计、开发以及测试工作，本平台已经顺利实现了既定的三端功能目标。

学生端已经形成了咨询师浏览与筛选、在线预约与取消、AI 心理助手流式对话、校园公告查阅以及咨询评价提交这五大功能模块；咨询师端已经形成了个人专业档案维护、预约审核（确认/拒绝）、会话记录管理以及案例笔记撰写这四大功能模块；学校管理端则已经形成了咨询师账号管理（创建/启停）、平台数据统计以及公告发布管理这三大功能模块。三端共同使用统一的后端 API 服务，并且借助 JWT 认证以及角色权限中间件来保障数据安全与访问隔离。

\subsection{技术创新点}

\textbf{SSE 流式对话技术的工程实践。}和在请求-响应模式下需要依靠轮询的方案相比，系统凭借 Server-Sent Events 技术建立起由服务端指向客户端的单向长连接，并且把 DeepSeek 大语言模型的流式输出实时转发到浏览器端，由此来实现逐词呈现的"打字机效果"。在这个过程当中，既涉及 Node.js 端对 \texttt{ReadableStream} 进行逐块读取以及 SSE 格式封装，也涉及前端对 \texttt{fetch} + \texttt{ReadableStream} 开展异步解析，还包括对多字节 UTF-8 字符跨块场景进行正确处理，所以整体上具有一定的工程复杂度。

\textbf{三端统一后端架构。}依靠模块化路由以及角色守卫中间件，三类用户可以共同使用同一个 Express 应用实例，这样一来就避免了多后端服务带来的部署复杂性以及数据同步问题。如果后续需要新增业务端点，只需在对应的路由文件里继续补充配置，就不会对其他端原有的业务逻辑造成影响，因此整体的扩展成本依然很低。

\textbf{三层 utf8mb4 编码加固方案。}针对 Node.js 与 MySQL 在连接过程当中较为常见的中文乱码问题，分别从连接配置层（\texttt{dialectOptions.charset}）、连接建立层（借助 \texttt{afterConnect} Hook 执行 \texttt{SET NAMES}）以及数据库表结构层这三个维度开展同步加固，这样一来就能够把由字符集不一致所引发的乱码隐患彻底消除，该方案对同类项目也有一定的参考价值。

\textbf{DATE 字段时区偏移的规避方案。}通过在 Sequelize 的 \texttt{dialectOptions} 里启用 \texttt{dateStrings: ['DATE']} 这一配置，让 mysql2 驱动在返回 \texttt{DATE} 类型字段时直接保留 \texttt{'YYYY-MM-DD'} 的字符串形式，这样就从驱动层把时区换算带来的日期偏移问题规避掉了，同时也不需要在应用层再额外加入补偿逻辑。

\subsection{系统价值}

从用户价值这一角度来看，平台把心理求助的起步环节从"预约排队等待数天"缩短到"打开页面就能对话"，7×24 小时 AI 助手的接入也因此明显降低了学生主动求助的门槛；凭借对预约信息以及会话记录开展数字化管理，咨询师在档案管理工作上的负担有所减轻，进而可以把更多精力放回专业服务本身；结构化的数据统计同时也为学校管理层提供了相对更客观的服务质量图景。

从技术价值这一角度来看，该系统相对完整地展示了一个前后端分离 Web 应用从数据库设计、API 规范以及认证鉴权一直到 AI 集成的整套全链路工程实践，所以对于同类高校心理健康信息化系统的建设来说，具有一定的参考意义。


\section{不足与展望}

\subsection{当前系统的不足}

坦率来说，鉴于开发周期以及能力边界方面的限制，现有系统在一些方面依然存在不足之处，因此还未达到生产级应用所要求的完备程度。

\textbf{功能方面。}系统目前还不支持用户自行去修改密码，用户在忘记密码之后也没有找回流程，只能依靠管理员来进行重置；头像上传的功能暂时仍未接入，所以界面的个性化呈现相对较弱；咨询师在收到新的预约时缺少系统推送提醒，只能借助手动刷新页面来查看，实时性有所不足；学生在结束会话之后也不会被系统主动引导去填写评价，这可能会使评价数据的收集不够完整。

\textbf{数据统计层面。}管理端目前已实现基于 ECharts 的多维图表统计，涵盖咨询师预约量分布、学生预约量排行及咨询话题占比，但统计视角仍局限于全量历史数据的静态汇总，尚未提供围绕时间维度的趋势分析，也缺少统计数据的导出功能，因此在精细化运营场景下仍有提升空间。

\textbf{部署层面。}系统目前还没有形成较为完整的容器化部署方案，环境配置工作仍然需要依靠手动来完成，因此在多服务器或者云端部署的场景当中，它的复用性以及可复制性都还显得不太充足。

\subsection{未来改进方向}

\textbf{实时消息推送。}后续可以考虑把 WebSocket 或服务端推送（Web Push API）这类机制引入进来，让咨询师在接收到新的预约时能够马上获得浏览器通知，进而对当前主要依靠手动刷新页面的被动体验进行改善。WebSocket 所具备的双向通信能力，也为后续去实现在线实时咨询这一文字聊天功能，提前打下了相应的技术基础。

\textbf{咨询师排班日历。}后续计划把当前以 JSON 格式展示的可预约时段，进一步升级为可视化的排班日历，让咨询师可以借助图形化方式开展每周排班的管理工作，并且在出现预约冲突时给出更为直观的提示，这样一来也能对操作体验进行提升。

\textbf{数据分析与导出。}现有图表统计已覆盖咨询师预约量、学生预约排行及话题分布等维度，后续可进一步增加以月为单位的预约趋势折线图与满意度评分走势图，帮助管理层识别周期性规律；同时支持将统计数据导出为 Excel 格式，方便开展周期性汇报与资料归档。

\textbf{移动端适配。}现有前端虽然已经完成了响应式布局框架的相关配置，不过在小屏幕设备上的交互细节方面，仍然还有进一步优化的空间。后续计划围绕 AI 对话以及预约流程等核心页面开展专项的移动端适配，或者以微信小程序作为载体来提供轻量版入口，从而把学生的使用门槛进一步降低。

\textbf{容器化部署。}后续计划提供一套基于 Docker Compose 的全栈部署方案，把前端静态服务、后端 API 服务以及 MySQL 数据库统一纳入编排流程，让整个系统的环境配置由手动操作转变为借助一条命令就能完成，进而对运维成本进行降低，同时也便于在不同高校之间复制推广。


\section*{本章小结}

本章围绕系统总结以及不足与展望这两个方面来展开，对本研究的整体推进过程进行了较为系统的回顾，同时也对后续可能的发展方向作出了前瞻性的梳理。系统总结部分对三端功能的整体落地情况进行了梳理，并且概括了 SSE 流式对话、三端统一架构、utf8mb4 三层加固以及 DATE 时区规避这四项技术创新，同时说明了系统在用户价值以及技术价值两个层面上的实践意义；不足分析部分也较为客观地指出，当前版本在功能完备性、数据统计深度以及部署便捷性这些方面仍然存在一定局限；展望部分则进一步提出了实时消息推送、排班日历、数据导出、移动端适配以及容器化部署这五个具备可操作性的扩展方向，为后续研究以及系统迭代提供了较为清晰的路线图。

